\documentclass[11pt]{scrartcl}
\usepackage[utf8]{inputenc}
\usepackage[sexy]{evanchen} 
\usepackage{graphicx}
\graphicspath{{8 Outsource K Math/Gambar/}}


\title{Modulo}
\author{Azzam and Akeyla}
\date{Maret 2023}

\begin{document}

\maketitle

\section{Modulo}
Misalkan $a,b,n$ adalah bilangan bulat dengan $n \neq 0$. Kita menyebut $a$ kongruen $b$ modulo $n$ apabila $n|a - b$. Kita menotasikan ini dengan $a \equiv b \pmod{n}$. 
\subsection{Properti}
\begin{enumerate}
\item $a \equiv a \pmod{n}$
\item Jika $a \equiv b \pmod{n}$ dan $b \equiv c \pmod{n}$, maka $a \equiv c \pmod{n}$
\item Jika $a \equiv b \pmod{n}$, maka $b \equiv a \pmod{n}$ 
\item Jika $a \equiv b \pmod{n}$, maka untuk sebarang bilangan bulat $k$, berlaku $ka \equiv kb \pmod{n}$
\item Jika $a \equiv b \pmod{n}$ dan $c \equiv d \pmod{n}$, maka untuk sebarang bilangan bulat $j,k$, berlaku $ja + kc \equiv jb + kd \pmod{n}$
\item Jika $a \equiv b \pmod{n}$ dan $c \equiv d \pmod{n}$, maka $ac \equiv bd \pmod{n}$
\item $a \equiv b \pmod{n_i}$ untuk $i = 1,2,3,..,k$ jika dan hanya jika $a \equiv b \pmod{ KPK(n_1,n_2,...,n_k)} $
\end{enumerate}

\subsection{Beberapa teorema yang berguna}
\begin{enumerate}
\item Untuk setiap bilangan asli $m,n$ yang relatif prima pasti terdapat suatu bilangan asli $x$ sedemikian hingga $mx \equiv 1$ (mod $n$) 
\item (Teorema Wilson) Untuk setiap bilangan prima $p$, berlaku $(p -1)! \equiv -1$ (mod $p$).
\item (Teorema Euler) Misalkan $\varphi (m)$ adalah banyak bilangan asli yang kurang dari $m$ dan relatif prima dengan $m$. Apabila faktorisasi prima dari $n$ adalah \[n = p_1^{\alpha_1}p_2^{\alpha_2}...p_k^{\alpha_k},\] maka nilai dari $\varphi(n)$ adalah \[n\left(1 - \dfrac{1}{p_1}\right)\left(1 - \dfrac{1}{p_2}\right)...\left(1 - \dfrac{1}{p_k}\right)\] Teorema Euler mengatkan bahwa untuk $a$ dan $m$ bilangan asli yang saling relatif prima, berlaku $a^{\varphi (m)} \equiv 1$ (mod $m$)
\item (Teorema Fermat Kecil/Fermat's Little Theorem) Untuk setiap bilangan prima $p$ dan bilangan asli $a$, berlaku $a^p \equiv a$ (mod $p$)
\end{enumerate}

\subsection{Contoh soal}
\begin{problem}
    Carilah sisa pembagian $2^{100}$ oleh $13$
\end{problem}

\begin{problem}
    (AIME 1983) Carilah sisa pembagian $6^{83} + 8^{83}$ oleh 49.    
\end{problem}

\subsection{Latihan Modulo}
\begin{problem}
    (Argentina MO 2021 Level 1) Mikha menulis barisan aritmetika yang dimulai dari $1$ dengan selisih $3$, yakni $$1, 4, 7, 10, 13, 16, ...$$ Kemudian, Fakhru menandai semua bilangan yang ditulis Mikha yang lebih dari $10$ dan kurang dari $100000$ dan seluruh digitnya sama. Tentukan banyak bilangan yang ditandai oleh Fakhru.
\end{problem}

\begin{problem}
    Tentukan dua digit terakhir dari $7^{707}$
\end{problem}

\begin{problem}
    Carilah digit terakhir dari $13^{14}$ dan $13^{14^{15}}$
\end{problem}

\begin{problem}
    (Canadian Math Olympiad Qualification Repechage 2016) Tentukan semua bilangan asli $n$ sehingga $11|(3^n + 4^n)$
\end{problem}

% \begin{problem}
%     Carilah dua digit terakhir dari $1032^{1032}$ dan $17^{17^{17}}$
% \end{problem}

\begin{problem}
    Tentukan sisa bagi saat $1\times 3\times 5\times ...\times 2023$ dibagi oleh $1000$.
\end{problem}
% \item Misalkan $p$ adalah bilangan prima ganjil. Untuk setiap bilangan $i = 1,2,..., p -1$, notasikan $r_i$ sebagai sisa pembagian ketika $i^p$ dibagi dengan $p^2$. Carilah nilai dari $$r_1 + r_2 + ... + r_{p -1}$$
% \item Misalkan $m,a,b$ adalah bilangan bulat dengan $a,b$ relatif prima dengan $m$. Bilangan asli $x$ dan $y$ adalah dua bilangan bulat sedemikian hingga $a^x \equiv b^x (mod m)$ dan $a^y \equiv b^y(mod m)$. Tunjukkan bahwa $$a^{FPB(x,y)} \equiv b^{FPB(x,y)}(mod m)$$
% \item Misalkan $a$ adalah bilangan ganjil. Buktikan bahwa $a^{2^n} + 2^{2^n}$ dan $a^{2^m} + 2^{2^m}$ relatif prima saat $m\neq n$.
% \item Jika $p$ adalah bilangan prima ganjil dan $a,b$ adalah dua bilangan asli yang relatif prima, maka tunjukkan bahwa $FPB(a+b, \dfrac{a^p + b^p}{a + b}) = 1$ atau $p$
% \item Misalkan $p \geq 7$ adalah bilangan prima. Tunjukkan bahwa bilangan $111...11$ (dengan $p$ buah angka 1) habis dibagi oleh $p$
% \item Misalkan $p$ adalah bilangan prima yang mempunyai bentuk $3k + 2$ yang membagi $a^2 + ab + b^2$ untuk sebuah bilangan bulat $a$ dan $b$. Buktikan bahwa $a$ dan $b$ keduanya habis dibagi oleh $p$.
% \item Buktikan bahwa untuk ssetiap bilangan genap positif $n$, berlaku $n^2 - 1$ membagi $2^{n!} - 1$
% \item Misalkan $x,a,b$ adalah bilangan asli. Tunjukkan bahwa $FPB(x^a - 1, x^b - 1) =x^{FPB(a,b)} - 1$.
% \item Misalkan $a$ dan $b>2$ adalah bilangan asli. Tunjukkan bahwa $2^a + 1$ tidak habis dibagi oleh $2^b - 1$
% \item Butikan bahwa untuk suatu bilangan prima $p$, terdapat suatu bilangan asli $x$ sedemikian hingga $x^2 + 1$ habis dibagi oleh $p$ jika dan hanya jika $p \equiv 1$ (mod 4).
% \item Misalkan $n > 1$ adalah bilangan ganjil. Tunjukkan bahwa $n$ tidak membagi $3^n + 1$
% \item Misalkan $n$ adalah bilangan asli dan $p$ bilangan prima sedemikian hingga $n$ membagi $p - 1$ dan $p$ membagi $n^3 - 1$. Tunjukkan bahwa $4p -3$ adalah bilangan kuadrat
% \item Misalkan $n>2$ adalah bilangan asli. Tunjukkan bahwa $$n^{n^{n^n}} - n^{n^n}$$ habis dibagi oleh 1989.
% \item Suatu titik letis $(x,y)$ di bidang kartesius disebut sebagai \textit{checkpoint} apabila $FPB(x,y) = 1$. Tunjukkan bahwa untuk setiap bilanga asli $n$ terdapat suatu persegi yang sisinya sejajar dengan sumbu-x dan sumbu-y sedemikian hingga seluruh titik letis yang ada di interior persegi merupakan \textit{checkpoint}.
% \item Apakah untuk setiap himpunan bilangan asli $\{a_1, a_2,...,,a_k\}$ terdapat suatu bilangan asli $b$ sedemikian hingga himpunan $\{ba_1,ba_2,...,ba_k\}$ memuat semua bilangan berpangkat sempurna?
% \item Misalkan $a$ adalah bilangan asli dan $p$ bilangan prima. Tunjukkan bahwa terdapat suatu bilangan asli $m$ yang memenuhi $$m^{m^m} \equiv a (mod p)$$
% \item Buktikan bahwa untuk setiap bilangan asli $n$, terdapat $n$ buah bilangan asli berurutan yang tidak ada satupun di antara $n$ bilangan tersebut yang merupakan bilangan berpangkat sempurna.
% \item Misalkan $a,b$ adalah dua bilangan asli sedemikian hingga $b^n + n$ adalah kelipatan dari $a^n + n$ untuk setiap bilangan asli $n$. Tunjukkan bahwa $a = b$.

\section{Bilangan Prima}
Suatu bilangan disebut sebagai bilangan prima apabila hanya ada tepat satu bilangan $d>1$ yang membagi $p$, yakni $d = p$. Bilangan asli yang $n>1$ disebut sebagai komposit jika $n$ bukan bilangan prima.
\subsection{Teorema Faktorisasi Prima}
Setiap bilangan asli $n > 1$ dapat dipresentasikan secara unik sebagai perkalian dari bilangan- bilangan prima.\\
Sebagai akibat dari teorema di atas, setiap biangan asli $n > 1$ dapat dinyatakan secara unik dalam bentuk $$n = p_1^{\alpha_1}p_2^{\alpha_2}...p_k^{\alpha_k}$$ dimana $p_1,p_2,...,p_k$ adalah bilangan prima berbeda dan $\alpha_1,\alpha_2,...,\alpha_k$ adalah bilangan asli.
\subsection{Beberapa Properti yang Sering Digunakan}
\begin{enumerate}
\item Banyak bilangan prima adalah tak hingga
\item Setiap bilangan prima yang lebih besar dari $3$ pasti berbentuk $6k \pm 1$ untuk suatu bilangan asli $k$
\item Jika $a,b$ adalah bilangan bulat dan $p$ bilangan prima sehingga $p|ab$, maka $p|a$ atau $p|b$ (atau keduanya)
\end{enumerate}

\subsection{Latihan Bilangan Prima}
\begin{problem}
    Tentukan semua bilangan asli $n$ sedemikian hingga $3n - 4, 4n - 5$, dan $5n - 3$ semuanya merupakan bilangan prima
\end{problem}

\begin{problem}
    (OSK 2011) Tentukan semua bilangan asli $p$ sedemikian hingga $p, p + 8$, dan $p + 16$ semuanya merupakan bilangan prima
\end{problem}

\begin{problem}
    Tentukan semua bilangan asli $a, b$ sedemikian hingga $a^4 + 4b^4$ merupakann bilangan prima.
\end{problem}

\begin{problem}
    (OSP 2010) Tentukan semua bilangan prima $p$ sehingga $p^2 + 7^3$ merupakan bilangan kubik.
\end{problem}

\begin{problem}
    (Argentina MO 2021 Level 1) Tentukan semua bilangan prima $p$ dan $q$ yang lebih dari $1$ dan kurang dari $100$ sedemikian hingga lima bilangan berikut: $$p + 6, p + 10, q + 4, q + 10, p + q + 1$$ semuanya merupakan bilangan prima
\end{problem}

% \subsection{Soal-Soal}
% \begin{enumerate}
% \item Buktikan bahwa bilangan prima $> 3$ selalu berbentuk $6k \pm 1$ untuk sutau bilangan asli $k$.
% \item Buktikan bahwa untuk $p>3$ prima, $24|p^2 - 1$
% \item Tunjukkan bahwa untuk $a,b$ bilangan asli,  $a^4 + 4b^4$ merupakan bilangan prima hanya saat $a = b = 1$
% \item Jika $p$ adalah bilangan prima ganjil dan $\dfrac{a}{b} = \dfrac{1}{1} + \dfrac{1}{2} + ... + \dfrac{1}{p - 1}$, buktikan bahwa $p|a$
% \item Buktikan bahwa jika $a^n + 1$ dan $a$ bilangan asli yang lebih besar dari 1, maka $a$ genap dan $n = 2^k$ untuk suatu bilangan cacah $k$. Bilangan prima dalam bentuk $2^{2^k} + 1$ disebut sebagai bilangan prima Fermat.
% \item Buktikan bahwa jika $a^n - 1$ dan $a$ bilangan asli yang lebih besar dari 1, maka $a = 2$ dan $n$ prima. Bilangan prima dalam bentuk $2^n - 1$ disebut sebagai bilangan prima Mersenne.
% \item Misalkan $p,q,r$ adalah bilangan prima yang memenuhi $p^2 + q^2 = r^n$ dengan $n$ adalah bilangan asli. Tunjukkan bahwa $n = 1$
% \end{enumerate}


% \section{Invers Modulo}
% Untuk $m,n$ dua bilangan asli yang relatif prima. Notasikan $m^{-1}$ (mod$n$) atau $\frac{1}{m}$ (mod$n$) sebagai nilai $x$(mod$n$) yang memenuhi $mx \equiv 1$ (mod$n$).

\section{Bilangan Berpangkat}
\subsection{Pemfaktoran dan Bilangan Berpangkat}
\begin{example}
    Bilangan cacah $n$ yang memenuhi $n^2 + 100$ merupakan bilangan kuadrat adalah ...
\end{example}

\begin{proof}[Solusi]
    Misalkan $n^2 + 100 = m^2$ untuk suatu bilangan asli $m$. Maka, akan diperoleh $m^2 - n^2 = 100$ yang dapat difaktorkan menjadi $(m - n)(m + n) = 100$. Sehingga, kita cukup mencari faktor dari $100$.
    \begin{table}[h!]
        \centering
        \begin{tabular}{c|c|c|c}
            \hline
            $m + n$ & $m - n$ & m & n  \\
            \hline 
            $100$ & $1$ & - & - \\
            $50$ & $2$ & $51$ & $49$ \\
            $25$ & $4$ & - & - \\
            $20$ & $5$ & - & - \\
            $10$ & $10$ & $10$ & $0$ \\
            $5$ & $20$ & - & - \\
            $4$ & $25$ & - & - \\
            $2$ & $50$ & - & - \\
            $1$ & $100$ & - & - \\
            \hline
        \end{tabular}
        \label{tab:my_label}
    \end{table}\\
    Dapat dilihat pada tabel, nilai $n$ yang memenuhi hanyalah $n = 0$ dan $n = 49$
\end{proof}

\subsection{Bilangan berpangkat dalam berbagai modulo}
Saat sebuah bilangan kuadrat dibagi dengan berbagai bilangan asli, tidak semua kemungkinan sisa modulo dapat diperoleh. Berikut adalah beberapa contohnya:
\begin{enumerate}
    \item Sisa bagi saat $x^2$ dibagi oleh $3$ adalah $0$ atau $1$
    \item Sisa bagi saat $x^2$ dibagi oleh $4$ adalah $0$ atau $1$
    \item Sisa bagi saat $x^2$ dibagi oleh $5$ adalah $0, 1,$ atau $4$
    \item Sisa bagi saat $x^2$ dibagi oleh $8$ adalah $0, 1,$ atau $4$
    \item Sisa bagi saat $x^2$ dibagi oleh $10$ adalah $0, 1, 4, 6,$ atau $9$
\end{enumerate}
Sifat sisa bagi yang tidak mencakup semua kemungkinan ini juga ditemui dalam pangkat yang lebih besar  seperti pada bilangan kubik. Sebagai contoh:
\begin{enumerate}
    \item Sisa bagi saat $x^3$ dibagi oleh $7$ adalah $0, 1,$ atau $6$
    \item Sisa bagi saat $x^4$ dibagi oleh $8$ adalah $0$ atau $1$
\end{enumerate}

\begin{example}
    Banyak bilangan asli $(a, b, c)$ yang memenuhi persamaan $a^2 + b^2 - 8c = 6$ adalah ...    
\end{example}

\begin{proof}[Solusi]
    Jawabannya adalah tidak ada. Perhatikan bahwa apabila ditinjau dalam modulo $8$, persamaan tersebut menjadi \[a^2 + b^2 \equiv \pmod{6}\]
    Namun, apabila ditinjau sisa bagi bilangan kuadrat saat dibagi oleh $8$, diperoleh bahwa sisa bagi yang memungkinkan adalah $0, 1,$ atau $4$. Penjumlahan dua sisa bagi manapun tidak akan dapat menghasilkan $6 \pmod{8}$. Sehingga, tidak ada bilangan $a, b, c$ yang memenuhi syarat soal.
\end{proof}

\subsection{Latihan Bilangan Berpangkat}
\begin{problem}
    Banyak bilangan di barisan \[11, 111, 1111, 11111, ..., 11...11\] dengan bilangan terakhir terdiri dari $2022$ digit $1$, yang merupakan bilangan kuadrat adalah ...
\end{problem}

\begin{problem}
    Carilah semua bilangan asli $m,n$ sedemikian hingga $6^m + 2^n + 2$ adalah bilangan kuadrat
\end{problem}

\begin{problem}
    Tentukan semua bilangan asli $n$ sehingga $n! + 5$ merupakan bilangan kubik
\end{problem}

% \begin{enumerate}
%     \item 
%     \item Tentukan banyak bilangan asli $x$ yang kurang dari 2020 sedemikian hingga $12x$ dan $20x$ keduanya merupakan bilangan kuadrat.
%     \item Tentukan semua bilangan asli $n$ sehingga $n^2 - 19n + 99$ merupakan bilangan kuadrat.
%     \item Carilah semua pasangan bilangan asli $(x,y)$ yang memenuhi $x^2 - y! = 2001$
% \end{enumerate}

\end{document}
